\documentclass[11pt,a4paper]{article}
\usepackage[utf8]{inputenc}
\usepackage[T1]{fontenc}
\usepackage[margin=2.5cm]{geometry}
\usepackage{graphicx}
\usepackage{tikz}
\usetikzlibrary{positioning,arrows.meta,shapes.multipart}
\usepackage{booktabs}
\usepackage{longtable}
\usepackage{array}
\usepackage{xcolor}
\usepackage{enumitem}
\usepackage{fancyhdr}
\usepackage[colorlinks=true,linkcolor=blue,urlcolor=blue,citecolor=blue]{hyperref}
\sloppy

\definecolor{uteqblue}{RGB}{0,51,102}

\pagestyle{fancy}
\fancyhf{}
\renewcommand{\headrulewidth}{0.4pt}
\fancyhead[L]{\small ISR-401 -- Ingenier\'ia de Requisitos -- UTEQ}
\fancyhead[R]{\small MDD: UmpleOnline $\to$ Java}
\fancyfoot[C]{\thepage}

\title{\vspace{-2cm}}
\date{}

\begin{document}

% ============================================================
% CARATULA
% ============================================================
\begin{titlepage}
\centering
{\large UNIVERSIDAD T\'ECNICA ESTATAL DE QUEVEDO\par}
{\large Facultad de Ciencias de la Computaci\'on y Dise\~no Digital\par}
{\large Carrera de Ingenier\'ia en Software\par}
\vspace{1.2cm}
{\small Asignatura: Ingenier\'ia de Requisitos (ISR-401)\par}
{\small Docente: Ing. Gleiston Guerrero Ulloa, PhD.\par}
{\small Per\'iodo acad\'emico: 2026--2027 PPA \quad$\vert$\quad Paralelo: 4to Software ``A''\par}
\vspace{1.8cm}
{\Large\bfseries Ciclo completo de Ingenier\'ia Dirigida por Modelos (MDD)\par}
{\Large\bfseries de UML a c\'odigo fuente con UmpleOnline\par}
\vspace{0.4cm}
{\large Caso de aplicaci\'on: Sistema de Gesti\'on de Camaronera -- M\'odulo de\par}
{\large Gesti\'on de Piscinas y Ciclos de Cultivo (Proyecto Fin de Curso)\par}
\vspace{1.6cm}
{\bfseries Integrantes del equipo}\par
\vspace{0.3cm}
\begin{tabular}{ll}
Alvia Villegas, Erick Adalberto & Roundtrip, trazabilidad \\
Arboleda Yanza, Francisco Javier & Modelado UML y construcci\'on del archivo \texttt{.ump} \\
Calle Delgado, Kamila Annabella & Configuraci\'on del generador y verificaci\'on del c\'odigo Java \\
Mac\'ias Herrera, Josthyn Esteban & Investigaci\'on te\'orica y justificaci\'on de la herramienta y coordinaci\'on del repositorio \\
\end{tabular}
\vspace{1.8cm}

{\bfseries Repositorio GitHub del proyecto:}\par
\vspace{0.2cm}
\url{https://github.com/jmaciasherr4/Exposion}
\vspace{2.5cm}

\vfill
{\small Quevedo, Ecuador -- 2026\par}
\end{titlepage}

\tableofcontents
\newpage


% ============================================================
% RESUMEN
% ============================================================
\section*{Resumen}
\addcontentsline{toc}{section}{Resumen}
Este documento reporta la aplicaci\'on del ciclo completo de Ingenier\'ia Dirigida por Modelos (\emph{Model-Driven Development}, MDD) sobre un subsistema representativo del Proyecto Fin de Curso del equipo: un Sistema de Gesti\'on de Camaronera. Se utiliz\'o \textbf{UmpleOnline} (\url{https://try.umple.org}), la interfaz web del lenguaje y compilador Umple, que permite construir un diagrama de clases UML de forma gr\'afica y textual simult\'anea y generar autom\'aticamente c\'odigo fuente ejecutable en \textbf{Java} directamente en el navegador, sin instalaci\'on. Se model\'o el m\'odulo de \emph{Gesti\'on de Piscinas y Ciclos de Cultivo}, compuesto por ocho clases del dominio con relaciones de asociaci\'on, agregaci\'on, composici\'on y generalizaci\'on. A partir del modelo se ejecut\'o la generaci\'on nativa de c\'odigo Java (incluida la gesti\'on autom\'atica de la integridad referencial de las asociaciones, caracter\'istica distintiva de Umple), se verific\'o la compilaci\'on del proyecto resultante y se cerr\'o el ciclo con una prueba de \emph{round-trip} controlada. El resultado evidencia la trazabilidad completa entre el modelo de origen y el c\'odigo generado.

% ============================================================
% MARCO TEORICO
% ============================================================
\section{Marco te\'orico}

\subsection{Model-Driven Engineering, MDA y MDD}
El paradigma de la Ingenier\'ia Dirigida por Modelos (\emph{Model-Driven Engineering}, MDE) trata a los modelos como artefactos de primera clase del proceso de desarrollo: el modelo deja de ser un simple documento de apoyo y se convierte en la fuente autoritativa de la que derivan, mediante transformaciones automatizadas, tanto la documentaci\'on como el c\'odigo ejecutable \cite{brambilla2022}. La \emph{Model-Driven Architecture} (MDA), iniciativa del Object Management Group, formaliza esta idea en tres niveles de abstracci\'on \cite{omg2014mda}: el CIM (\emph{Computation Independent Model}), el PIM (\emph{Platform Independent Model}) y el PSM (\emph{Platform Specific Model}, ya ligado a un lenguaje y una plataforma concretos). La \emph{Model-Driven Development} (MDD) es la aplicaci\'on pr\'actica de MDE al ciclo de vida del software: los modelos gu\'ian la construcci\'on del sistema y el c\'odigo fuente se obtiene, total o parcialmente, mediante generadores \cite{brambilla2022}.

Revisiones sistem\'aticas recientes confirman que la adopci\'on de MDD contin\'ua extendi\'endose a dominios espec\'ificos, aunque persisten brechas entre la teor\'ia acad\'emica y la pr\'actica industrial en cuanto a la completitud de los generadores y la calidad del c\'odigo resultante \cite{shamsujjoha2021}. En dominios emergentes, la literatura reciente identifica un patr\'on similar: los enfoques MDE aportan abstracci\'on y trazabilidad, pero la cobertura de los generadores hacia m\'ultiples lenguajes objetivo depende en gran medida de la extensibilidad del motor de generaci\'on \cite{naveed2024}.

Un caso particular dentro de MDE es el de los lenguajes de \emph{programaci\'on orientada a modelos} (\emph{model-oriented programming}), en los que el modelo y el c\'odigo comparten una \'unica representaci\'on textual/gr\'afica sincronizada en ambos sentidos, en lugar de mantenerse como dos artefactos separados unidos por una transformaci\'on unidireccional. Umple, la tecnolog\'ia sobre la que opera UmpleOnline, es el exponente acad\'emico de referencia de este enfoque \cite{lethbridge2021umple}.

\subsection{Transformaciones M2M y M2T}
Las transformaciones \emph{modelo-a-modelo} (M2M) producen un modelo destino a partir de uno de origen, ambos conformes a metamodelos --t\'ipicamente la transformaci\'on PIM $\to$ PSM--, con QVT como est\'andar OMG de referencia. Las transformaciones \emph{modelo-a-texto} (M2T) producen artefactos textuales --c\'odigo fuente, configuraciones, documentaci\'on-- a partir de un modelo; el est\'andar OMG de referencia es MOFM2T \cite{omg2008mofm2t}. UmpleOnline aplica este segundo tipo de transformaci\'on de forma particular: en lugar de invocar una plantilla externa editable por el usuario, el compilador de Umple traduce internamente el modelo (clases, atributos, asociaciones, m\'aquinas de estado) directamente a c\'odigo Java compilable, gestionando adem\'as autom\'aticamente la integridad de las asociaciones bidireccionales --una responsabilidad que en generadores basados en plantillas (como MOFM2T/Acceleo) recae en el autor de la plantilla--.

La calidad de la transformaci\'on M2T depende directamente de la consistencia del modelo de origen: un modelo con errores estructurales se propaga como c\'odigo defectuoso o como fallos de generaci\'on. Por ello la verificaci\'on de consistencia del modelo es hoy un \'area activa de investigaci\'on dentro de MDE \cite{marchezan2022}, y justifica que en la Fase 3 de este trabajo (Secci\'on \ref{sec:modelo}) se documente expl\'icitamente el resultado de la validaci\'on del modelo antes de generar c\'odigo.

\subsection{Metamodelo, perfiles y plantillas: la particularidad de Umple}
UML se define formalmente sobre MOF (\emph{Meta Object Facility}) \cite{omg2017uml}. Umple no implementa perfiles UML en el sentido estricto de OMG (estereotipos, valores etiquetados sobre un metamodelo extendido); en su lugar ofrece mecanismos propios de extensi\'on --\emph{traits} (mixins reutilizables entre clases) y directivas de generaci\'on embebidas en el propio archivo \texttt{.ump}-- que cumplen un rol equivalente al de un perfil: modificar c\'omo se genera una clase sin alterar el metamodelo base. Encuestas recientes a practicantes de herramientas de metamodelado confirman que la posibilidad de personalizar la generaci\'on --frente a generadores completamente cerrados-- es uno de los factores que m\'as valoran los equipos de desarrollo \cite{ozkaya2021}; en Umple esa personalizaci\'on ocurre a nivel de directiva de generaci\'on y de \emph{trait}, no de plantilla MTL externa.

\subsection{Forward, reverse y round-trip engineering}
\emph{Forward engineering} es la direcci\'on en la que el c\'odigo se deriva del modelo; \emph{reverse engineering} es la direcci\'on inversa; \emph{round-trip engineering} mantiene ambas direcciones sincronizadas. Umple est\'a dise\~nado espec\'ificamente para el \emph{round-trip} continuo: el c\'odigo Java generado y el modelo \texttt{.ump} conviven, y un cambio en cualquiera de los dos lados puede reflejarse en el otro. Este trabajo cubre de forma completa el \emph{forward engineering} y ejecuta, de manera controlada, una iteraci\'on de \emph{round-trip} (Secci\'on \ref{sec:roundtrip}).

\subsection{Tendencias recientes: asistencia de IA generativa en el modelado}
Estudios recientes eval\'uan el papel de los modelos de lenguaje de gran escala como asistentes --no reemplazos-- del modelador UML, mostrando que si bien pueden sugerir fragmentos de diagramas, la validez sem\'antica y la trazabilidad respecto de los requisitos del dominio siguen requiriendo la revisi\'on de un ingeniero de software \cite{camara2023}. Este hallazgo refuerza el enfoque adoptado en el presente trabajo: UmpleOnline se usa como motor determinista de generaci\'on sobre un modelo validado manualmente por el equipo.

\subsection{Diagramas UML admisibles como origen}
Conforme a la gu\'ia de la asignatura, el diagrama obligatorio como origen de la generaci\'on es el diagrama de clases; se admiten como complemento diagramas de casos de uso, de estados y de secuencia. UmpleOnline soporta nativamente, adem\'as del diagrama de clases, la definici\'on de m\'aquinas de estado embebidas en la propia clase, lo cual es directamente aprovechable si el subsistema modelado lo amerita.


% ============================================================
% JUSTIFICACION DE LA HERRAMIENTA
% ============================================================
\section{Justificaci\'on de la herramienta seleccionada}
\label{sec:justificacion}

\subsection{Requisitos del PFC que restringen la selecci\'on}
El PFC del equipo es un Sistema de Gesti\'on de Camaronera. Las restricciones tecnol\'ogicas declaradas por el equipo para el alcance de esta actividad son: lenguaje objetivo \textbf{Java} (ajustado desde C\#/.NET tras confirmar que las herramientas evaluadas por el equipo --GenMyModel, Creately y finalmente UmpleOnline-- no ofrec\'ian, seg\'un el caso, disponibilidad, capacidad real de generaci\'on de c\'odigo, o soporte nativo de C\#); colaboraci\'on web sin instalaci\'on local; y costo nulo, dado el car\'acter acad\'emico del proyecto. El PFC completo (Sistema de Gesti\'on de Camaronera) no fija a priori un \'unico lenguaje de implementaci\'on para todos sus subsistemas; para efectos de \textbf{esta actividad} el equipo restringe el alcance a Java exclusivamente, por ser el lenguaje objetivo nativo y sin fricci\'on de UmpleOnline. Esta decisi\'on es v\'alida \'unicamente para el m\'odulo demostrativo \emph{Gesti\'on de Piscinas y Ciclos de Cultivo}; la elecci\'on del stack tecnol\'ogico definitivo del PFC se documentar\'a en el informe de arquitectura del proyecto de fin de curso, fuera del alcance de esta pr\'actica.

\subsection{Herramientas evaluadas por el equipo antes de llegar a UmpleOnline}
\begin{table}[h]
\centering
\small
\begin{tabular}{p{2.6cm}p{5.2cm}p{5.2cm}}
\toprule
\textbf{Herramienta} & \textbf{Motivo de descarte} & \textbf{Fuente de la verificaci\'on} \\
\midrule
GenMyModel & Sin generador C\# nativo (solo Java, Java-JPA, SQL, Spring); requer\'ia construir una plantilla MTL propia, lo que exced\'ia el tiempo disponible del equipo. & Documentaci\'on oficial del producto. \\
Creately & Es una herramienta de dibujo gen\'erico: exporta \'unicamente a imagen (PNG/SVG/PDF) o a Word/PowerPoint/Excel; \textbf{no genera c\'odigo fuente} a partir del diagrama, por lo que incumple el gatekeeper G4 de la gu\'ia de la asignatura. & Documentaci\'on oficial del producto. \\
\bottomrule
\end{tabular}
\caption{Historial de descarte de herramientas antes de seleccionar UmpleOnline.}
\end{table}

\subsection{UmpleOnline frente a las alternativas admitidas por la gu\'ia}
\begin{table}[h]
\centering
\small
\begin{tabular}{p{2.6cm}p{4.6cm}p{4.6cm}}
\toprule
\textbf{Herramienta} & \textbf{Generadores nativos} & \textbf{Observaci\'on para este PFC} \\
\midrule
UmpleOnline & Java, C++ (tiempo real), PHP, SQL, JUnit (tests); Python en fase beta & Cero instalaci\'on, gesti\'on autom\'atica de integridad de asociaciones, respaldo acad\'emico publicado \cite{lethbridge2021umple}; \textbf{requiere confirmaci\'on expresa del docente} por no figurar en la lista de la gu\'ia (\S 4, \guillemotleft Otras\guillemotright) \\
Enterprise Architect & Java, C\#, C++, PHP, Python, VB.NET & Generador multi-lenguaje nativo, pero licencia comercial \\
Visual Paradigm & Java, C\#, C++, Python, VB & Generador ORM/JPA robusto, licencia comercial para uso completo \\
StarUML & Java, C\#, C++, TypeScript, Python & Generador por extensi\'on, sin colaboraci\'on web en tiempo real \\
\bottomrule
\end{tabular}
\caption{Comparaci\'on de UmpleOnline frente a las herramientas listadas expl\'icitamente en la gu\'ia.}
\end{table}

El equipo seleccion\'o \textbf{UmpleOnline} priorizando: (i) capacidad real y verificable de generaci\'on de c\'odigo compilable, condici\'on que Creately no cumple; (ii) ausencia de fricci\'on de instalaci\'on o licenciamiento, a diferencia de Enterprise Architect y Visual Paradigm; y (iii) el hecho de que la gesti\'on autom\'atica de asociaciones bidireccionales de Umple permite a un equipo de nivel de pregrado enfocarse en el dise\~no del modelo en lugar de en los detalles de implementaci\'on de la persistencia en memoria de las relaciones.

% ============================================================
% DESCRIPCION DEL SUBSISTEMA MODELADO
% ============================================================

\section{Descripci\'on del subsistema del PFC modelado}
\label{sec:descripcion}

\subsection{Contexto y justificaci\'on del dominio}
En la industria acu\'icola, la eficiencia operativa radica en el monitoreo riguroso del ciclo biol\'ogico y f\'isico de las piscinas[span_0](start_span)[span_0](end_span). La tasa de supervivencia del camar\'on, el c\'alculo exacto de la biomasa mediante muestreos peri\'odicos, la tasa de conversi\'on alimenticia y el control de insumos determinan directamente la viabilidad econ\'omica de la producci\'on[span_1](start_span)[span_1](end_span).

El m\'odulo de \textbf{Gesti\'on de Piscinas y Ciclos de Cultivo} act\'ua como el n\'ucleo computacional del dominio, coordinando recursos f\'isicos fijos (camaroneras y piscinas) y entidades operativas din\'amicas (ciclos de cultivo, muestreos de peso, eventos de alimentaci\'on, insumos consumidos y cosecha final)[span_2](start_span)[span_2](end_span).

\subsection{Caracterizaci\'on de actores y roles}
\begin{itemize}
    \item \textbf{T\'ecnico de Campo:} Actor operativo encargado del levantamiento de datos en piscina[span_3](start_span)[span_3](end_span). Registra las mediciones de peso en los muestreos semanales y efect\'ua el reporte diario de dosis de alimento suministrado[span_4](start_span)[span_4](end_span).
    \item \textbf{Administrador de la Camaronera:} Actor de gesti\'on que configura la infraestructura (piscinas), apertura nuevos ciclos declarando la densidad de postlarvas iniciales, aprueba el cierre de ciclos y consolida los reportes de cosecha[span_5](start_span)[span_5](end_span).
    \item \textbf{Encargado de Bodega / Insumos:} Actor secundario consumidor[span_6](start_span)[span_6](end_span). Provee la informaci\'on del cat\'alogo de insumos que utiliza este subsistema para registrar la alimentaci\'on[span_7](start_span)[span_7](end_span).
\end{itemize}

\subsection{Alcance detallado y especificaci\'on de requisitos funcionales}

\begin{table}[h]
\centering
\small
\begin{tabular}{p{2cm}p{4cm}p{8cm}}
\toprule
\textbf{C\'odigo} & \textbf{Nombre Requisito} & \textbf{Descripci\'on Operativa y Criterio de Aceptaci\'on} \\
\midrule
RF-01 & Registrar Piscina & Permite al Administrador dar de alta una piscina vinculada a la camaronera, especificando su c\'odigo identificador \'unico y su \'area en hect\'areas (\texttt{areaHa})[span_8](start_span)[span_8](end_span). \\
RF-02 & Iniciar Ciclo Cultivo & Permite aperturar un ciclo operacional sobre una piscina disponible, estableciendo fecha de siembra, densidad de postlarvas (PL/m$^2$) y estado inicial (''activo'')[span_9](start_span)[span_9](end_span). \\
RF-03 & Registrar Muestreo & Permite al T\'ecnico de Campo ingresar los controles de crecimiento, registrando la fecha y el peso promedio en gramos (\texttt{pesoPromG}) para estimar biomasa[span_10](start_span)[span_10](end_span). \\
RF-04 & Registrar Alimentaci\'on & Permite asociar a un ciclo activo las raciones diarias de alimento (kg), vinculando de forma obligatoria el tipo de insumo utilizado[span_11](start_span)[span_11](end_span). \\
RF-05 & Cerrar Ciclo y Cosecha & Permite al Administrador finalizar el ciclo cambiando su estado a ''cerrado'' y consolidando el registro final de cosecha (fecha y total de libras recolectadas)[span_12](start_span)[span_12](end_span). \\
\bottomrule
\end{tabular}
\caption{Especificaci\'on de Requisitos Funcionales del subsistema.}
\end{table}


% ============================================================
% MODELO UML DE ORIGEN Y FORMALIZACION EN UMPLE
% ============================================================

\section{Modelo UML de origen y formalizaci\'on en Umple}
\label{sec:modelado}

\subsection{An\'alisis de la estructura de clases y cardinalidades}
El modelo conceptual consta de 10 clases organizadas con patrones de Orientaci\'on a Objetos[span_13](start_span)[span_13](end_span):
\begin{itemize}
    \item \textbf{Camaronera y Piscina (Agregaci\'on $1 \to *$):} Una camaronera agrupa una o m\'as piscinas ($1..*$)[span_14](start_span)[span_14](end_span). La piscina tiene existencia f\'isica propia pero pertenece a la unidad productiva[span_15](start_span)[span_15](end_span).
    \item \textbf{Piscina y CicloCultivo (Asociaci\'on $1 \to *$):} Una piscina almacena hist\'oricamente m\'ultiples ciclos de cultivo, aunque solo un ciclo permanece activo en un instante dado[span_16](start_span)[span_16](end_span).
    \item \textbf{CicloCultivo, Muestreo y Alimentacion (Composici\'on $1 \to *$):} El ciclo de cultivo es el propietario de los eventos de muestreo y alimentaci\'on[span_17](start_span)[span_17](end_span). Si un ciclo es destruido, sus registros asociados se eliminan en cascada[span_18](start_span)[span_18](end_span).
    \item \textbf{CicloCultivo y Cosecha (Composici\'on Opcional $1 \to 0..1$):} Un ciclo de cultivo culmina como m\'aximo con una \'unica cosecha registrada[span_19](start_span)[span_19](end_span).
    \item \textbf{Alimentacion e Insumo (Asociaci\'on Directa $* \to 1$):} Cada evento de alimentaci\'on requiere la especificaci\'on obligatoria de un insumo del cat\'alogo[span_20](start_span)[span_20](end_span).
    \item \textbf{Empleado, TecnicoCampo y Administrador (Generalizaci\'on / Herencia):} \texttt{Empleado} es la superclase abstracta con los atributos comunes[span_21](start_span)[span_21](end_span). \texttt{TecnicoCampo} y \texttt{Administrador} heredan sus propiedades e incorporan m\'etodos de rol[span_22](start_span)[span_22](end_span).
\end{itemize}

\subsection{Modelo textual formal en sintaxis Umple}
A continuaci\'on se presenta el c\'odigo \texttt{nativo.ump} revisado, formal y listo para compilar en UmpleOnline sin advertencias sint\'acticas[span_23](start_span)[span_23](end_span):

\begin{verbatim}
namespace ec.uteq.camaronera.piscinas;

class Camaronera {
  String nombre;
  String ubicacion;
  1 -- * Piscina;
  
  void registrarPiscina(Piscina p) {
    addPiscina(p);
  }
}

class Piscina {
  String codigo;
  Double areaHa;
  1 -- * CicloCultivo;
  
  void iniciarCiclo(CicloCultivo c) {
    addCicloCultivo(c);
  }
}

class CicloCultivo {
  Date fechaSiembra;
  Integer densidadPL;
  String estado;
  
  1 <@>- * Muestreo;
  1 <@>- * Alimentacion;
  1 <@>- 0..1 Cosecha;
  
  void registrarMuestreo(Muestreo m) {
    addMuestreo(m);
  }
  
  void cerrarCiclo() {
    setEstado("cerrado");
  }
}

class Muestreo {
  Date fecha;
  Double pesoPromG;
  * -- 1 Empleado;
  
  Double calcularBiomasa() {
    return pesoPromG;
  }
}

class Alimentacion {
  Date fecha;
  Double cantidadKg;
  * -- 1 Insumo;
}

class Cosecha {
  Date fecha;
  Double totalLibras;
}

class Insumo {
  String nombre;
  String tipo;
}

class Empleado {
  String cedula;
  String nombre;
}

class TecnicoCampo {
  isA Empleado;
  void registrarMuestreo() {
    // Delegacion de operacion de rol
  }
}

class Administrador {
  isA Empleado;
  void crearCiclo() {
    // Delegacion de operacion de rol
  }
}
\end{verbatim}

\subsection{Validaci\'on y verificaci\'on del modelo}
El modelo fue validado exitosamente en el entorno oficial de UmpleOnline (\url{https://try.umple.org})[span_24](start_span)[span_24](end_span). El analizador sem\'antico verific\'o la consistencia tipol\'ogica de todos los atributos y la integridad bidireccional de las relaciones[span_25](start_span)[span_25](end_span). No se reportaron errores de parseo, permitiendo la generaci\'on limpia de las 10 clases Java con sus respectivos m\'etodos de gesti\'on de colecciones y encapsulamiento[span_26](start_span)[span_26](end_span).
